% ============================================================
%  EDM 2026 --- Speaker notes for the 7.5-minute talk
%  Paper: Graph-Based vs Transition-Based Dependency Parsing
%         for Uzbek (A Cross-Architecture Study under
%         Low-Resource Conditions)
%  Authors: S. Matlatipov, M. Ismoilova, M. Aripov
%
%  Optimised for printing on a black & white printer:
%    - 12pt serif body, 1.4 line spacing
%    - bold cues only (no colours), clear slide separators
%    - generous left margin for handwritten annotations
%    - cumulative timing budget printed next to each slide
%
%  Build with:  pdflatex presentation_speaker_notes.tex
% ============================================================
\documentclass[12pt,a4paper]{article}

\usepackage[utf8]{inputenc}
\usepackage[T1]{fontenc}
\usepackage{lmodern}
\usepackage[margin=2.2cm,left=3.2cm,right=2.2cm]{geometry}
\usepackage{setspace}
\usepackage{enumitem}
\usepackage{titlesec}
\usepackage{fancyhdr}
\usepackage{microtype}
\usepackage{parskip}
\usepackage{amsmath,amssymb}

\setstretch{1.40}
\setlist[itemize]{leftmargin=1.3em,topsep=2pt,itemsep=2pt}

\titleformat{\section}
   {\normalfont\Large\bfseries\filright}
   {Slide \thesection.}{0.6em}{}
\titlespacing*{\section}{0pt}{1.2em}{0.4em}

\pagestyle{fancy}
\fancyhf{}
\fancyhead[L]{\small EDM 2026 --- Speaker notes (7.5 min)}
\fancyhead[R]{\small S.~Matlatipov}
\fancyfoot[C]{\thepage}
\renewcommand{\headrulewidth}{0.4pt}

\newcommand{\timing}[2]{%
   \par\noindent\textbf{\small [#1\,$\to$\,#2\,min]}\par\vspace{0.2em}}
\newcommand{\say}[1]{\textbf{#1}}
\newcommand{\breathe}{\par\vspace{0.4em}\textbf{[pause]}\par}

\begin{document}

% ------------------------------------------------------------
%  COVER PAGE
% ------------------------------------------------------------
\begin{center}
   {\LARGE\bfseries Speaker notes --- 7.5 minutes}\\[0.6em]
   {\large Graph-Based vs Transition-Based\\Dependency Parsing for Uzbek}\\[0.3em]
   {\normalsize A Cross-Architecture Study under Low-Resource Conditions}\\[0.6em]
   {\normalsize International Conference on Educational Data Mining}\\
   {\normalsize \textbf{EDM 2026}}\\[1.2em]
   \textit{Sanatbek Matlatipov, Maftuna Ismoilova, Mersaid Aripov}\\
   National University of Uzbekistan, Tashkent\\
   \texttt{s.matlatipov@nuu.uz}
\end{center}

\vspace{1em}
\noindent\textbf{How to use these notes.}
Each numbered section is one slide. The bracketed times
\textbf{[start\,$\to$\,end]} are cumulative running time against a strict
7.5-minute budget. Bold phrases are the words to land on; \textbf{[pause]}
is a deliberate beat. The prose is a spoken script --- deliver it
conversationally, not read verbatim.

\vspace{0.5em}
\noindent\textbf{Discipline for a short slot.} Average $\sim$30\,s per
slide and keep momentum. Do \emph{not} read tables cell-by-cell --- point
to the one figure that carries the argument. Protect time on
\textbf{Slide~7 (subword fusion)} and \textbf{Slide~10 (the 16.70 LAS
gap)}; everything else can move briskly.

\vspace{0.5em}
\noindent\textbf{The single claim to leave behind:} \emph{for a
morphologically rich, genuinely low-resource language, the choice of
parsing \underline{architecture} matters more than tagging accuracy ---
and cross-treebank data is the most powerful lever we have.}
\textbf{Numbers to land:} \textbf{LAS 63.81 vs 47.11}, a
\textbf{16.70-point} gap; and cross-treebank \textbf{+9.62 LAS} for the
graph-based parser.

\vfill\hrule\vspace{0.5em}
\noindent\textit{Print double-sided on A4. The wide left margin is for handwritten cues.}

\newpage

% ------------------------------------------------------------
%  SLIDE 1 --- Title
% ------------------------------------------------------------
\section*{Slide 1 --- Title}
\timing{0:00}{0:20}

Good morning, and thank you to the EDM organisers and to all of you for
being here.

My name is \say{Sanatbek Matlatipov}, from the National University of
Uzbekistan, presenting joint work with Maftuna Ismoilova and Mersaid
Aripov. Our question is deceptively simple, and for Uzbek it has never
been answered: \say{when you have almost no training data, which family
of neural dependency parser should you actually reach for --- graph-based
or transition-based?} And, just as importantly, \say{why}.

% ------------------------------------------------------------
%  SLIDE 2 --- Why is parsing Uzbek hard?
% ------------------------------------------------------------
\section*{Slide 2 --- Why is parsing Uzbek hard?}
\timing{0:20}{1:05}

Three properties make this hard. First, Uzbek is \say{agglutinative}: a
single word stacks an ordered chain of suffixes encoding case, tense,
person, even evidentiality. Second, it has \say{flexible
subject--object--verb} word order with frequent center-embedding, plus
null subjects and no articles --- so surface cues English parsers rely on
simply are not there. Third, and decisively, it is \say{low-resource}:
the entire Universal Dependencies data for Uzbek is just \textbf{1,184
sentences}.

\breathe

The box on the right is the intuition that drives the whole talk. Take
\textit{bolalarning} --- ``of the children.'' It decomposes into the stem
\textit{bola}, the plural \textit{-lar}, and the genitive
\textbf{\textit{-ning}}. That final suffix is not decoration --- it
\emph{is} the syntactic relation. So in Uzbek, \say{if you lose the last
suffix, you lose the parse}. Hold onto that; it returns on Slide~7.

% ------------------------------------------------------------
%  SLIDE 3 --- Contributions
% ------------------------------------------------------------
\section*{Slide 3 --- Contributions}
\timing{1:05}{1:35}

We make three contributions. First, the \say{first head-to-head
comparison} of a graph-based parser, Stanza, against a transition-based
parser, spaCy, for Uzbek --- and critically, \say{both wired to the same
TahrirchiBERT encoder}, so the architecture is the \emph{only} variable.
Second, we show that cross-treebank augmentation helps the two
architectures \say{very differently} --- a finding, not a footnote.
Third, we give \say{per-relation and per-feature error analyses} that
explain \emph{where} each architecture wins. Data, code, and checkpoints
are all public.

% ------------------------------------------------------------
%  SLIDE 4 --- Two UD treebanks for Uzbek
% ------------------------------------------------------------
\section*{Slide 4 --- Two UD treebanks for Uzbek}
\timing{1:35}{2:05}

We work with two treebanks. \say{UzUDT} --- 684 sentences of fiction and
academic prose, hand-annotated by six people on INCEpTION with
inter-annotator agreement above 95\%. And \say{Uzbek-UT} --- 500
sentences of news and fiction, semi-automatic with full manual
correction. Neither shipped with a standard split, so we re-split both
consistently. Merged, that is \textbf{781 training sentences} --- which
roughly \say{doubles} the data over UzUDT alone. Keep that doubling in
mind for the augmentation result later.

% ------------------------------------------------------------
%  SLIDE 5 --- Why the treebanks are complementary
% ------------------------------------------------------------
\section*{Slide 5 --- Why the treebanks are complementary}
\timing{2:05}{2:35}

But this is not just \emph{more} data --- it is \say{broader} data, and
that distinction is the heart of the augmentation story. The two
treebanks cover genuinely different syntax. UT is news-heavy: far more
proper nouns and \say{ten times} the light-verb compounds. UzUDT is
literary: \say{four times} the adverbial clauses, plus possessor
agreement and evidentiality that UT simply does not annotate. So they
fill each other's blind spots --- the union covers more of the grammar
than either alone.

% ------------------------------------------------------------
%  SLIDE 6 --- Two neural parsing architectures
% ------------------------------------------------------------
\section*{Slide 6 --- Two neural parsing architectures}
\timing{2:35}{3:15}

Same encoder, two fundamentally different decoders. \say{Stanza} is
graph-based: a two-stage pipeline ending in a Dozat--Manning biaffine
parser that scores \emph{every} possible head--dependent pair and then
decodes the best tree globally. \say{spaCy} is transition-based: a single
jointly-trained model that builds the tree incrementally with arc-eager
shift--reduce moves. They even differ in how they collapse BERT
sub-words into tokens --- Stanza takes the last sub-word, spaCy averages.
That sub-word choice turns out to matter, so let me take it next.

% ------------------------------------------------------------
%  SLIDE 7 --- Subword to token fusion (KEY)
% ------------------------------------------------------------
\section*{Slide 7 --- Subword-to-token fusion}
\timing{3:15}{4:05}

This is the slide to slow down on. BERT emits \say{WordPieces}, but
Universal Dependencies annotates \say{whole tokens}, and the two do not
align. So we must pool several sub-word vectors into one per token --- and
\emph{how} we pool is a real modelling decision.

Look at the table. In every example the \emph{last} piece ---
\textit{\#\#ning}, \textit{\#\#ni}, \textit{\#\#ib} --- carries the
grammatically decisive morpheme: genitive, accusative, converb. This is
not a coincidence; it is a typological fact about a suffixing language.

\breathe

So \say{mean pooling} --- spaCy's default --- averages every piece and
\emph{dilutes} exactly that suffix signal. \say{Last-sub-word} selection
keeps it intact. The payoff is clean: in a controlled ablation,
last-sub-word buys \textbf{+2.64 LAS} with \say{zero} extra parameters ---
and the gain shows up in \emph{parsing}, not tagging, precisely because
attachment in Uzbek is case-driven.

% ------------------------------------------------------------
%  SLIDE 8 --- Experimental setup
% ------------------------------------------------------------
\section*{Slide 8 --- Experimental setup}
\timing{4:05}{4:30}

Briefly, the protocol. Seven models: five Stanza configurations, two
spaCy. The suffix \say{point-one} means UzUDT only, \say{point-two} means
merged. We deliberately use each framework's \say{documented defaults} ---
this is a fair out-of-the-box comparison, not a tuning contest, which is
how a practitioner actually meets these tools. One GPU, standard UD
metrics: UPOS, XPOS, morphological features, and unlabeled and labeled
attachment scores.

% ------------------------------------------------------------
%  SLIDE 9 --- Results within each framework
% ------------------------------------------------------------
\section*{Slide 9 --- Results within each framework}
\timing{4:30}{5:05}

Two quick reads before the comparison. \say{Within Stanza}, on the left:
contextual BERT beats static FastText on tagging, and the highlighted row
--- \say{E2.2}, BERT with last-sub-word on merged data --- is our best
configuration at \textbf{LAS 63.81}. \say{Within spaCy}, on the right:
again, merged data wins. And notice the pattern that holds in \emph{both}
frameworks --- the single biggest jump comes from merging, up to eleven
LAS points. That sets up the two questions the next slides answer.

% ------------------------------------------------------------
%  SLIDE 10 --- The central comparison (CLIMAX)
% ------------------------------------------------------------
\section*{Slide 10 --- The central comparison}
\timing{5:05}{5:55}

This is the core result of the paper. Same encoder, same merged data ---
\say{only the architecture differs}.

On tagging, spaCy is \emph{better}: UPOS of 89 against 85. But look at the
last row. On \say{labeled parsing}, Stanza scores \textbf{63.81} against
spaCy's \textbf{47.11} --- a \say{16.70-point LAS gap}. That is among the
largest architecture gaps reported in the recent parsing literature.

\breathe

And here is the punchline I want you to take away: Stanza wins parsing
\emph{despite} \say{losing} on tagging. So for Uzbek, raw POS accuracy is
\emph{not} what drives parse quality --- the \say{exhaustive biaffine
search} is, because a rich case system and long-distance dependencies
reward scoring all head--dependent pairs globally rather than committing
greedily, left to right.

% ------------------------------------------------------------
%  SLIDE 11 --- Data augmentation is architecture-dependent
% ------------------------------------------------------------
\section*{Slide 11 --- Data augmentation is architecture-dependent}
\timing{5:55}{6:25}

Now the second finding: merging helps both architectures, but \say{not
in the same place}. For tagging the gains are nearly identical --- about
+2.6 UPOS either way. For \emph{parsing}, they diverge sharply: Stanza
gains \textbf{+9.62 LAS}, spaCy only +1.76 --- the graph-based parser is
the one that actually exploits the structural diversity we added.
Symmetrically, spaCy's \emph{morphology} leaps \textbf{+14.93}, because
its joint objective feeds parser gradients back into the morphologiser.
Architecture decides \emph{what} extra data buys you.

% ------------------------------------------------------------
%  SLIDE 12 --- Per-relation error analysis
% ------------------------------------------------------------
\section*{Slide 12 --- Per-relation error analysis}
\timing{6:25}{6:55}

Where do the errors concentrate? The frequency plot on the left shows a
heavily skewed relation distribution. The per-relation scores on the
right tell the story: the transition-based parser is \say{strong on
local, short relations} --- \textit{nummod}, \textit{root} --- but
\say{weak on the long-distance ones}: coordination, \textit{conj}, drops
to 18.9, and noun modification, \textit{nmod}, suffers from the genitive
\textit{-ning} ambiguity we saw on Slide~2. These are exactly the
structurally hard cases that exhaustive graph scoring handles better.

% ------------------------------------------------------------
%  SLIDE 13 --- Qualitative parses
% ------------------------------------------------------------
\section*{Slide 13 --- Qualitative parses}
\timing{6:55}{7:15}

Two real outputs from our best system to make this concrete. On the left,
a \say{perfect parse} --- 100\% --- with correct genitive--possessive
chaining and a converb \textit{xcomp}: the diagnostic Uzbek
constructions, all right. On the right, a \say{systematic failure} --- the
light-verb construction \textit{davom etmoq}, ``to continue.'' Gold
\textit{compound:lvc} is in blue; the model's wrong arcs are dashed in
red. Rare multi-word predicates like this are precisely what a tiny
treebank under-covers --- and precisely where more, broader data helps.

% ------------------------------------------------------------
%  SLIDE 14 --- Conclusion
% ------------------------------------------------------------
\section*{Slide 14 --- Conclusion}
\timing{7:15}{7:30}

To conclude. We find a clear \say{architecture-dependent trade-off}:
transition-based parsing tags better, but graph-based parsing parses far
better --- \textbf{+16.70 LAS}. Cross-treebank augmentation is the
\say{single most impactful lever}, and graph-based parsers exploit it
best. So our practical recommendation for Uzbek and languages like it is
unambiguous: \say{a graph-based parser, a monolingual BERT, and
last-sub-word fusion}, whenever parse quality is the goal. Next steps are
multi-seed significance testing, broader frameworks and encoders, and
LLM-based data augmentation.

% ------------------------------------------------------------
%  SLIDE 15 --- Thank you
% ------------------------------------------------------------
\section*{Slide 15 --- Thank you}
\timing{7:30}{--}

Thank you. The treebanks, the code, and all trained checkpoints are
released --- on Hugging Face and GitHub --- and I would be glad to take
your questions. \textit{(Contact: \texttt{s.matlatipov@nuu.uz}.)}

% ------------------------------------------------------------
%  Q & A prep --- printed as a back page for the speaker only
% ------------------------------------------------------------
\newpage
\section*{Likely Q\&A --- quick cues}

\noindent\textbf{Q. Is the comparison fair if the hyper-parameters differ?}\\
\textbf{Yes, by design.} We use each framework's documented defaults to
reflect the genuine out-of-the-box practitioner experience. On a corpus
this small, tuning rarely flips an architectural ranking, and a
16.70-point LAS gap dwarfs any plausible tuning effect.

\vspace{0.4em}
\noindent\textbf{Q. Why does Stanza freeze BERT while spaCy fine-tunes it?}\\
That reflects each framework's \textbf{native training paradigm}, not a
choice we imposed: Stanza trains a parser head over frozen features;
spaCy fine-tunes the transformer jointly. We report both as practitioners
would actually run them.

\vspace{0.4em}
\noindent\textbf{Q. These are single runs --- is the 16.70 LAS gap real?}\\
\textbf{Yes.} A paired-bootstrap on the 325-sentence test set gives a
95\% confidence interval of about $\pm$2.4 LAS per model --- far below the
gap --- and the ranking is monotone across all five Stanza runs.
Multi-seed testing is our immediate next step.

\vspace{0.4em}
\noindent\textbf{Q. Why last-sub-word rather than mean or attention pooling?}\\
It is \textbf{motivated by Uzbek morphology}: the rightmost sub-word
carries the outermost, syntactically decisive suffix. Empirically it is
+2.64 LAS over mean pooling with no extra parameters, and it should
transfer to other suffixing Turkic languages.

\vspace{0.4em}
\noindent\textbf{Q. Why not UDPipe 2.0, Trankit, or XLM-R / UzbBERT?}\\
This paper is a focused \textbf{graph-vs-transition} study, not a
multi-framework benchmark. Broader frameworks and alternative encoders
are explicit future work.

\vspace{0.4em}
\noindent\textbf{Q. Does merging help because of size or because of genre?}\\
\textbf{Both.} It nearly doubles the data \emph{and} the treebanks are
genre-complementary --- UT contributes proper nouns and light-verb
compounds, UzUDT contributes adverbial clauses and possessor/evidential
features.

\vspace{0.4em}
\noindent\textbf{Q. How does this connect to educational data mining?}\\
Robust Uzbek parsing is foundational infrastructure for downstream
educational NLP --- automated writing feedback, readability assessment,
and content analysis --- for a language with 37 million speakers and very
little tooling.

\end{document}
